\documentclass[11pt,letterpaper]{article}

\usepackage[margin=1in]{geometry}
\usepackage{amsmath,amssymb,amsthm}
\usepackage{bm}
\usepackage{xcolor}
\usepackage{tikz}
\usetikzlibrary{arrows.meta}
\usepackage[colorlinks=true,linkcolor=blue!60!black,urlcolor=blue!60!black,citecolor=blue!50!black]{hyperref}
\usepackage{cleveref}
\usepackage{enumitem}

\newcommand{\Jpm}{J_{\pm}}
\newcommand{\Jzz}{J_{zz}}
\newcommand{\bvarphi}{\bm{\varphi}}
\newcommand{\bw}{\mathbf{w}}
\newcommand{\bn}{\mathbf{n}}
\newcommand{\bN}{\mathbf{N}}
\newcommand{\br}{\mathbf{r}}
\newcommand{\bL}{\mathbf{L}}
\newcommand{\brho}{\bm{\rho}}
\newcommand{\bdelta}{\bm{\delta}}
\newcommand{\Hperp}{H_{\perp}}
\newcommand{\Heff}{H_{\mathrm{eff}}}
\newcommand{\gfour}{g_{4}}
\newcommand{\ghex}{g_{\mathrm{hex}}}
\newcommand{\Z}{\mathbb{Z}}
\newcommand{\twoinout}{2\text{-in-}2\text{-out}}

\theoremstyle{plain}
\newtheorem{theorem}{Theorem}
\newtheorem{lemma}{Lemma}
\newtheorem{corollary}{Corollary}
\newtheorem{proposition}{Proposition}
\theoremstyle{definition}
\newtheorem{definition}{Definition}
\newtheorem{remark}{Remark}

\title{Twist averaging is not enough:\\
virtual-path windings and the exact zero-transport projector\\
for the quantum spin ice gauge sector}
\author{}
\date{\today}

\begin{document}
\maketitle

\begin{abstract}
This note is the theory backbone of the accompanying program: turning a
small periodic quantum spin ice (QSI) cluster into a quantitative
laboratory for the emergent U(1) gauge sector by removing the
torus-winding processes that have no infinite-lattice analogue. An
earlier version of this document proved a \emph{Winding-Projection
Theorem} asserting that averaging over twisted boundary conditions
removes those processes exactly, for every closed-system observable.
\textbf{That theorem is false, and we retract it here.} Its error is a
single, instructive one: it assigned each emergent ring exchange the
holonomy of its \emph{loop}, whereas the twist phase of a perturbative
process is the product of the phases of the bond terms the process
\emph{actually uses} -- a property of virtual \emph{paths}, not loops.
We prove the correct structural result, the \emph{Matched-Pair
Factorization} (\cref{thm:matched}): a ring exchange over a cycle
$\mathcal C$ is generated by the several perfect matchings of
$\mathcal C$ into dimer moves, and its twisted coupling is the sum of
$e^{-i\bN\cdot\bvarphi}$ over matchings, with $\sum\bN=\bw_{\mathcal C}$
the loop winding. Two corollaries then pin down when averaging works. (A)~The operator
average is \emph{gauge-dependent}: in the boundary-bond gauge every
torus-winding ring exchange has a zero-wrap matching, so the average
retains one half of it; but in the translationally-invariant (smooth)
gauge the same average removes it entirely -- because there it becomes
the \emph{zero-transport projector} onto processes with net charge
transport $\bdelta=\brho+\bN\cdot\bL=\bm0$ (the gauge-invariant label a
uniform emergent vector potential couples to). (B)~The \emph{observable}
average -- what one computes in practice, averaging a specific $C(T)$ or
$S(q,\omega)$ -- fails \emph{gauge-invariantly}, keeping every even power
of the spurious coupling. So twist averaging removes the artifact if and
only if it is done in the smooth gauge at the operator level, at which
point it \emph{is} the zero-transport projector; the boundary gauge
halves it and any observable average keeps the full spurious scale. The
projector, being a statement about processes rather than spectra, cleans
not only $H$ but every ice-manifold response operator. A tight-binding toy solvable by hand
exhibits all three statements. We close with the specialization to QSI
and the role of the projector in the $\pi$-flux-beyond-QMC campaign.
\end{abstract}

\tableofcontents

%=====================================================================
\section{Plain-language summary}
\label{sec:nomath}

We simulate a magnet by cutting a small block out of the crystal and
gluing opposite faces together, so a spin leaving one face re-enters the
opposite one. The glued block is a three-dimensional torus, and
diagonalizing the spin Hamiltonian on it is \emph{exact diagonalization
on a periodic cluster}.

The gluing creates fake shortcuts: closed loops of spin flips that exist
only because the faces are identified, with no counterpart on the
infinite crystal. In QSI these fake loops are especially damaging
because they are \emph{short} -- they enter the low-energy (gauge)
dynamics at a lower order in the coupling than the genuine ring
exchange, so they can dominate exactly the quantity one wants.

The natural idea is to thread an Aharonov--Bohm flux $\bvarphi$ through
each wrap-around direction, a ``twist,'' and average the results over
$\bvarphi$. A loop that wraps picks up a phase; average over the dial
and the phase seems to average to zero. \textbf{This intuition is
wrong,} and the reason is the point of this note. An emergent ring
exchange is not a single connected walk around the loop; it is a product
of local moves, and the same net loop can be produced by several
distinct arrangements of those moves. One of those arrangements always
uses only interior (non-wrapping) moves -- so it carries \emph{no} twist
phase, survives any average, and leaves the fake loop half-alive. Worse,
in practice one averages a specific number like the specific heat, which
is a nonlinear function of the energies; averaging it cancels even less.

The cure is to look not at which loop a process draws, but at how much
electric charge it \emph{transports} across the cluster. A genuine local
ring exchange transports nothing; a fake wrapping loop necessarily
transports charge from one side to the other. Keeping only the
zero-transport processes deletes every fake loop exactly, at every order,
while leaving every genuine one untouched -- and it does so for every
physical quantity at once, because it is a rule about the processes
themselves, not about any one observable. The rest of this note makes
each italicized claim a theorem.

%=====================================================================
\section{Setup: the cluster torus and the ice manifold}
\label{sec:setup}

The XXZ pyrochlore Hamiltonian in the local frames is
\begin{equation}
H=\Jzz\sum_{\langle ij\rangle}S^z_iS^z_j
-\Jpm\sum_{\langle ij\rangle}\bigl(S^+_iS^-_j+S^-_iS^+_j\bigr),
\qquad |\Jpm|\ll\Jzz .
\label{eq:H}
\end{equation}
The cluster is a set of $N$ sites with periodic identifications
generated by three translation vectors collected as the rows of a
matrix $\bL\in\Z^{3\times3}$; the position space is the flat torus
$T^3=\mathbb R^3/\bL\Z^3$, with first homology $H_1(T^3,\Z)=\Z^3$. For a
nearest-neighbour bond $(i,j)$ let $\bn_{ij}=\mathrm{round}
[(\br_j-\br_i)\bL^{-1}]\in\Z^3$ be its \emph{wrap vector}: $\bn_{ij}=\bm
0$ for an interior bond, nonzero for a boundary-crossing one.

For $|\Jpm|\ll\Jzz$ the low-energy manifold is the space of
$\twoinout$ (ice-rule) configurations, $\dim=N_{\rm ice}$; degenerate
perturbation theory in the transverse term produces an effective ring
Hamiltonian $\Heff$ acting on it. On the infinite lattice the shortest
ice-preserving loop is the six-site hexagon, giving the standard photon
coupling $\ghex=12|\Jpm|^3/\Jzz^2$. On the torus the identifications
also admit shorter ice-preserving cycles that wind; the dominant ones
are four-site loops, second order in $\Jpm$, coupling
$\gfour=4\Jpm^2/\Jzz\gg\ghex$.

%=====================================================================
\section{A twist is a flat U(1) connection; its holonomy}
\label{sec:connection}

A twist $\bvarphi\in[0,2\pi)^3$ is inserted by attaching a phase to each
transverse hop according to the bond it uses,
\begin{equation}
\Hperp(\bvarphi)=-\Jpm\sum_{\langle ij\rangle}
\bigl(S^+_iS^-_j\,e^{-i\bn_{ij}\cdot\bvarphi}
+S^-_iS^+_j\,e^{+i\bn_{ij}\cdot\bvarphi}\bigr),
\label{eq:Htwist}
\end{equation}
in the \emph{minimal (boundary-bond) gauge}, where only wrapping bonds
carry phases. The Ising term is untouched, $S^z_{\rm tot}$ is conserved
at every $\bvarphi$, and $\bvarphi=\bm 0$ recovers \cref{eq:H}. Because
the connection is flat, its only gauge-invariant content is holonomy:
for a closed sequence of hops with total wrap
$\bw=\sum_a\bn_{i_a i_{a+1}}\in\Z^3$ the accumulated phase is the Wilson
loop $e^{-i\bw\cdot\bvarphi}$, a function of the homology class alone.

The tempting -- and false -- inference is that this makes each emergent
ring exchange carry $e^{-i\bw\cdot\bvarphi}$. The next section shows why
it does not: a ring exchange is not a single closed hop sequence.

%=====================================================================
\section{The engine: operators are sums over virtual paths}
\label{sec:pathexp}

Every quantity computed from $\Hperp(\bvarphi)$ -- matrix elements of
$\Heff$, the resolvent $(z-\Hperp)^{-1}$, the heat kernel
$e^{-\beta\Hperp}$ -- is, by expanding in $\Hperp$, a sum over sequences
of elementary hops, each sequence weighted by the product of its hop
amplitudes. Under \eqref{eq:Htwist} each hop $(i\!\to\!j)$ contributes a
phase $e^{-i\bn_{ij}\cdot\bvarphi}$. Hence a process that starts and ends
in the same configuration and uses hops of total wrap $\bN$ carries the
phase $e^{-i\bN\cdot\bvarphi}$.

The crucial point is that an \emph{operator} on the ice manifold -- an
emergent ring exchange -- is not one such sequence. The ring-exchange
operator over an ice-preserving cycle $\mathcal C=(a b c d\cdots)$,
\begin{equation}
O_{\mathcal C}=\bigl(S^+_aS^-_bS^+_cS^-_d\cdots\bigr),
\end{equation}
is a product of commuting local dimer moves, and it is generated in
perturbation theory by \emph{every} factorization of $\mathcal C$ into
disjoint nearest-neighbour dimer moves -- the \emph{perfect matchings}
of the cycle. Different matchings use different bonds, hence different
total wraps $\bN$. This is the entire content of the correction to
\cref{sec:connection}.

%=====================================================================
\section{The Matched-Pair Factorization Theorem}
\label{sec:master}

\begin{theorem}[Matched-Pair Factorization]
\label{thm:matched}
Let $\mathcal C$ be an ice-preserving simple cycle of even length
$2m$ on the cluster, with winding $\bw_{\mathcal C}$. In the minimal
gauge \eqref{eq:Htwist}, at the leading perturbative order at which
$O_{\mathcal C}$ appears, its coefficient in $\Heff(\bvarphi)$ is
\begin{equation}
c_{\mathcal C}(\bvarphi)
= -\,\frac{g_{\mathcal C}}{M}\sum_{\mu=1}^{M}
e^{-i\bN_\mu\cdot\bvarphi},
\qquad
\bN_\mu=\!\!\sum_{(ij)\in\text{matching }\mu}\!\!\bn_{ij},
\label{eq:matched}
\end{equation}
where the sum runs over the $M$ perfect matchings of $\mathcal C$ into
disjoint dimer moves, $g_{\mathcal C}>0$ is the (real, gauge-fixed)
bare coupling, and the matching wraps obey $\bN_1+\bN_2=\bw_{\mathcal C}$
for a two-matching loop. Only the individual matching wraps $\bN_\mu$
enter; the loop winding constrains only their sum. In
particular, the gauge-invariant magnitude for a two-matching loop
($M=2$) is
\begin{equation}
|c_{\mathcal C}(\bvarphi)|
=\frac{g_{\mathcal C}}{2}\bigl|1+e^{-i(\bN_2-\bN_1)\cdot\bvarphi}\bigr|
=g_{\mathcal C}\,\bigl|\cos\tfrac{\bw_{\mathcal C}\cdot\bvarphi}{2}\bigr|,
\label{eq:mag}
\end{equation}
using $\bN_2-\bN_1=\pm\bw_{\mathcal C}$.
\end{theorem}

\begin{proof}
At the leading order, $O_{\mathcal C}$ is produced only by processes
that flip each of the $2m$ spins of $\mathcal C$ exactly once, with no
other flips; these are exactly the products of $m$ disjoint dimer moves
covering $\mathcal C$, i.e.\ its perfect matchings. A cycle of even
length has $M\ge2$ matchings (for a $4$-cycle, $M=2$: $\{ab,cd\}$ and
$\{bc,da\}$). Each matching $\mu$ is a product of $m$ commuting hops; in
the minimal gauge its phase is the product $e^{-i\bN_\mu\cdot\bvarphi}$
with $\bN_\mu$ the sum of the constituent bond wraps. The real amplitudes
of all matchings coincide (equal denominators and combinatorial weight
at leading order), giving the prefactor $g_{\mathcal C}/M$. Finally,
$\bN_1+\bN_2$ for the two matchings of a $4$-cycle traverses each bond of
$\mathcal C$ once, i.e.\ equals the loop wrap $\bw_{\mathcal C}$; hence
$\bN_2-\bN_1\equiv\bw_{\mathcal C}\pmod2$ and \eqref{eq:mag} follows.
\end{proof}

\begin{remark}[Where the old proof failed]
The retracted Winding-Projection Theorem assumed $M=1$: that
$O_{\mathcal C}$ carries the single phase $e^{-i\bw_{\mathcal
C}\cdot\bvarphi}$. Equation~\eqref{eq:matched} shows the operator lives
in \emph{several} twist-Fourier sectors $\bN_\mu$ at once; the holonomy
$e^{-i\bw\cdot\bvarphi}$ appears only as the \emph{relative} phase of the
matchings, \eqref{eq:mag}. Everything below is a consequence.
\end{remark}

%=====================================================================
\section{Corollary A: twist-averaging the Hamiltonian keeps half}
\label{sec:cor-half}

\begin{corollary}[Half-survival, boundary-bond gauge]
\label{cor:half}
Work in the minimal (boundary-bond) gauge \eqref{eq:Htwist}, where only
wrapping bonds carry twist phases. Suppose $\mathcal C$ is a torus-winding
ice loop possessing a matching $\mu_0$ with $\bN_{\mu_0}=\bm 0$ (all its
dimer moves interior). Then
for \emph{any} averaging measure $d\nu(\bvarphi)$ on the twist torus --
the continuum $\int d^3\varphi/(2\pi)^3$, any regular grid, or the eight
corners $\{0,\pi\}^3$ --
\begin{equation}
\overline{c_{\mathcal C}}
=\int c_{\mathcal C}(\bvarphi)\,d\nu(\bvarphi)
= -\frac{g_{\mathcal C}}{M}\Bigl(1+\textstyle\sum_{\mu\ne\mu_0}
\widehat{\nu}(\bN_\mu)\Bigr),
\end{equation}
where $\widehat\nu(\bN)=\int e^{-i\bN\cdot\bvarphi}d\nu$. The $\mu_0$
term $-g_{\mathcal C}/M$ survives untouched, and for a two-matching loop
this is exactly one half of the bare coupling. Twist averaging of
$\Heff$ therefore cannot annihilate $O_{\mathcal C}$; at best it removes
the $M-1$ wrapping matchings.
\end{corollary}

\begin{proof}
$\widehat\nu(\bm0)=1$ for any normalized measure, so the $\bN_{\mu_0}=\bm
0$ term is invariant under averaging. The other terms are multiplied by
$\widehat\nu(\bN_\mu)$, of modulus $\le1$.
\end{proof}

\begin{proposition}[The hypothesis always holds on the production clusters]
\label{prop:census}
On the $1{\times}1{\times}1$ cubic ($16$-site) and $2{\times}2{\times}2$
FCC ($32$-site) pyrochlore clusters, every ice-preserving winding
four-loop and every wrapping hexagon has a matching with
$\bN_{\mu_0}=\bm0$, whereas every contractible hexagon has \emph{all}
matchings with $\bN_\mu=\bm0$.
\end{proposition}

This is verified by exhaustive enumeration (path-resolved
Schrieffer--Wolff, tracking $\bN$ per virtual path): on the cubic cluster
all $36$ four-loops and all $48$ wrapping hexagons retain $\tfrac12$
their coupling under the eight-corner operator average, while all $16$
contractible hexagons are preserved in full; the FCC cluster gives the
same with $24$, $96$, $32$. Thus twist averaging leaves the artifact at
the $\gfour/2$ scale -- a fact confirmed spectrally: the corner-averaged
$16$-site low-$T$ peak stalls at $0.018\,\Jzz\approx\gfour/2$, never
reaching $\ghex=0.012\,\Jzz$. A direct check on a single winding
$4$-loop [sites $(0,2,8,10)$, $\bw=(-1,0,-1)$] makes the mechanism
concrete: its exact ring coefficient obeys
$|c_{\mathcal C}(\bvarphi)|=\gfour\,|\cos(\bw\cdot\bvarphi/2)|$ to all
digits along $\bw$, dropping to \emph{zero} at $\bw\cdot\bvarphi=\pi$ --
a single homology phase $e^{-i\bw\cdot\bvarphi}$ (constant magnitude)
cannot vanish, whereas two interfering matchings do.

\begin{remark}[The operator average is gauge-dependent; the smooth gauge
is the projector]
\label{rem:gauge}
\cref{cor:half} is a statement about the \emph{boundary-bond} gauge, and
the ``one half'' is gauge-dependent -- as it must be, since the twist
average of an \emph{operator} averages a gauge-dependent family
$\Heff(\bvarphi)$. Choose instead the translationally-invariant
(``smooth'') gauge, a uniform vector potential $\mathbf A=\bL^{-1}\bvarphi$
with $\theta_{ij}=\mathbf A\cdot\mathbf d_{ij}$ on \emph{every} bond
(interior included). Then a process carries $e^{-i\mathbf A\cdot\bdelta}$
with $\bdelta=\brho+\bN\cdot\bL$ its charge transport
[\cref{eq:delta}], and averaging over the extended period projects onto
$\bdelta\equiv\bm0$: this is exactly the zero-transport projector of
\cref{thm:transport}. It removes the winding loop \emph{entirely}, not by
half. The three schemes give, for the loop above,
\begin{equation}
\underbrace{-\gfour}_{\text{bare }\bvarphi=0}
\;\longrightarrow\;
\underbrace{-\tfrac{\gfour}{2}}_{\text{boundary-gauge op.\ avg.}}
\;\longrightarrow\;
\underbrace{\;0\;}_{\text{smooth-gauge op.\ avg.\ }=\text{ projector}},
\end{equation}
while a contractible hexagon is preserved ($+\ghex$) by all three. Thus
twist averaging \emph{does} remove the artifact -- but only in the
translationally-invariant gauge and only at the operator level, at which
point it \emph{is} the projector of \cref{sec:transport}. The failure of
naive twist averaging is therefore twofold: the boundary gauge halves
rather than removes (\cref{cor:half}), and any \emph{observable} average
keeps even more (\cref{cor:obs} below), the latter gauge-invariantly.
\end{remark}

\begin{remark}[Why a winding operator can be twist-blind: bonds vs.\
tetrahedra]
\label{rem:bondtet}
It is natural to object that a winding operator, being non-existent under
open boundaries, ``ought'' to be caught by the twist. The resolution is
that the twist is a connection on \emph{bonds}, whereas the ice-rule
closure of a process is a property of \emph{tetrahedra}, and the two need
not coincide. For the loop above, the zero-wrap matching $\mu_0$ uses the
interior bonds $\{(0,2),(8,10)\}$ (neither crosses the boundary, so
$\bN_{\mu_0}=\bm0$ and it carries no twist phase), yet its virtual spinon
pair can only be reabsorbed on the tetrahedron $\{0,5,10,15\}$, whose
defining pair $(10,0)$ \emph{is} a wrapping bond. Under open boundaries
that tetrahedron is absent and the process fails to return to the ice
manifold -- so the operator does require periodicity -- but the twist,
living on bonds, never sees the interior matching. What does see it is
the charge \emph{transport} $\bdelta=\brho+\bN\cdot\bL$ of
\cref{sec:transport}: here $\bN_{\mu_0}=\bm0$ but the operator dipole
$\brho=\br_0+\br_8-\br_2-\br_{10}=(-\tfrac12,0,-\tfrac12)\neq\bm0$, so
$\bdelta\neq\bm0$ and the zero-transport projector removes exactly the
matching that twist averaging cannot.
\end{remark}

\begin{corollary}[Corner switch-off, not sign-flip]
\label{cor:corner}
At a corner $\bvarphi=\pi\mathbf p$, $\mathbf p\in\{0,1\}^3$,
\eqref{eq:mag} gives $|c_{\mathcal C}|=g_{\mathcal C}$ if
$\bw_{\mathcal C}\cdot\mathbf p$ is even and $0$ if odd: a winding loop's
coupling is switched fully \emph{off} at half the corners, not
sign-reversed. The eight-corner average is therefore a
loop-deactivation histogram, not a cancellation.
\end{corollary}

%=====================================================================
\section{Corollary B: averaging observables is weaker still}
\label{sec:cor-obs}

In practice one does not average $\Heff$; one averages an
\emph{observable} -- the specific heat $C(T,\bvarphi)$, the free energy,
a structure factor -- and these are nonlinear in $\Heff$.

\begin{corollary}[Even powers survive -- gauge-invariantly]
\label{cor:obs}
Write the spurious coupling of a single winding loop as
$g(\bvarphi)=g_{\mathcal C}\cos(\bw\cdot\bvarphi/2)$
[\cref{eq:mag}] and let $\mathcal O(g)$ be any closed-system observable,
analytic in $g$ near the physical value. Its eight-corner (or continuum)
average retains every \emph{even} power of $g$:
\begin{equation}
\overline{\mathcal O}
=\sum_{k\ge0}\mathcal O_{2k}\,\overline{g^{2k}},
\qquad \overline{g^{2k}}=g_{\mathcal C}^{2k}\,\overline{\cos^{2k}}\neq0,
\end{equation}
because $g(\bvarphi)$ is even about each corner. Odd powers, and only
they, can cancel. Since a winding loop enters observables already at
$O(g^2)$ -- e.g.\ a gap $2\sqrt{\varepsilon^2+g^2}$ about a genuine scale
$\varepsilon\sim\ghex$ -- and $g_{\mathcal C}=\gfour\gg\ghex$, the
surviving even-power contamination dominates the physical scale.
\end{corollary}

Thus the observable average is strictly weaker than the operator average
of \cref{sec:cor-half}: on $16$ sites the low-$T$ peak moves only from
$0.037\,\Jzz$ (bare) to $0.031\,\Jzz$ (corner-averaged observable),
against $0.018\,\Jzz$ for the corner-averaged operator -- neither
reaching $\ghex$.

%=====================================================================
\section{The zero-transport projector}
\label{sec:transport}

The correct invariant is not homology but transport. Define, for a
process taking ice configuration $c$ to $c'$, the \emph{net charge
transport}
\begin{equation}
\bdelta \;=\; \brho \;+\; \bN\cdot\bL,
\qquad
\brho=\sum_{i\,\text{raised}}\br_i-\sum_{j\,\text{lowered}}\br_j ,
\label{eq:delta}
\end{equation}
where $\brho$ is the dipole of the flipped-site set (the operator's
intrinsic displacement) and $\bN\cdot\bL$ the contribution of any wraps.
$\bdelta$ is the physical displacement of the transported $S^z$
(boson) density and is gauge invariant: it is exactly what a
\emph{uniform} emergent vector potential couples to. In the smooth gauge
$\theta_{ij}=\mathbf A\cdot\mathbf d_{ij}$ with $\bL\mathbf
A=\bvarphi$, every process carries $e^{-i\mathbf A\cdot\bdelta}$, and
averaging over the extended twist period is projection onto $\bdelta=\bm
0$.

\begin{theorem}[Zero-transport projector]
\label{thm:transport}
The projector
\begin{equation}
\Heff^{\bdelta}=\!\!\sum_{\text{processes}:\ \bdelta=\bm0}\!\!(\text{that process})
\end{equation}
annihilates every torus-winding ice ring exchange, at every order in
$\Jpm$, and preserves every contractible one, provided the operator
dipole is too short to bridge a cluster period ($|\brho|<\min_{\bw\ne\bm
0}|\bw\cdot\bL|$), which holds for the four-loops and hexagons of the
production clusters.
\end{theorem}

\begin{proof}
\emph{Winding loops die.} A matching of a ring loop transports the
flipped bosons by at most its diameter, $|\brho|\le$ (loop extent);
for a $4$-cycle this is $\le2d_{\rm NN}=a/\sqrt2$, for a hexagon
$\le3d_{\rm NN}$. A winding loop has $\bw\ne\bm0$, so its two matchings'
transports differ by $\bw\cdot\bL$, with
$|\bw\cdot\bL|\ge a$ (cubic) or $\sqrt2 a$ (FCC): since $|\brho|<
|\bw\cdot\bL|$, no matching can reach $\bdelta=\bm0$. Higher-order
dressings shift $\brho$ only by full lattice vectors, which the
short operator dipole cannot cancel; hence $\bdelta\ne\bm0$ for a
winding loop at every order.

\emph{Contractible loops survive.} A contractible ice loop has
$\bw=\bm0$ and, being a genuine bulk plaquette (a regular $[111]$
hexagon), has vanishing alternating displacement sum, $\brho=\bm0$, for
each matching. Thus $\bdelta=\bm0$ and the process is kept in full.
\end{proof}

\begin{remark}[Operator, not spectrum]
\cref{thm:transport} is a rule on processes, so it cleans \emph{every}
ice-manifold operator simultaneously -- $\Heff$ and every induced
response operator alike. This is what upgrades the projector from a
spectral fix to a definition of the bulk gauge sector of a torus, and is
the basis of the response program of the main paper.
\end{remark}

\begin{remark}[One honest bias]
The projector also deletes any \emph{contractible} process carrying a
nonzero transport dipole ($\brho\ne\bm0$), possible for
non-centrosymmetric loops at order $\Jpm^4$ and beyond. Within the
nearest-neighbour XXZ model these are $O(\Jpm^4/\Jzz^3)$, parametrically
below $\ghex$; we flag rather than hide them.
\end{remark}

%=====================================================================
\section{A toy model solvable by hand}
\label{sec:toy}

Take the two-dimensional space spanned by a single ice loop's two
ice configurations $\{|{+}\rangle,|{-}\rangle\}$, connected by a winding
$4$-loop with two matchings, and coupled to a genuine (contractible)
process of strength $\varepsilon$. By \cref{thm:matched} the winding
block at twist $\varphi$ (one relevant component) is
\begin{equation}
\Heff(\varphi)=
\begin{pmatrix} 0 & \varepsilon+\tfrac{\gfour}{2}(1+e^{-i\varphi})\\[2pt]
\text{h.c.} & 0\end{pmatrix},
\end{equation}
with $\bN_1=0$ (interior matching), $\bN_2=1$ (wrapping matching), so
the winding coupling is $g(\varphi)=\tfrac{\gfour}{2}(1+e^{-i\varphi})$.

\emph{Operator average} (\cref{cor:half}):
$\overline{g}=\tfrac{\gfour}{2}\bigl(1+\overline{e^{-i\varphi}}\bigr)
=\tfrac{\gfour}{2}$ for any average with $\overline{e^{-i\varphi}}=0$
(continuum, or the two ``corners'' $\varphi\in\{0,\pi\}$). Half survives.

\emph{Observable average} (\cref{cor:obs}): the gap is
$\Delta(\varphi)=2|\varepsilon+g(\varphi)|$, and
$\overline{\Delta}$ over $\varphi\in\{0,\pi\}$ is
$|\varepsilon+\gfour|+|\varepsilon|$ -- the $\gfour$ term never cancels.
Even less is removed than by the operator average.

\emph{Zero-transport projection} (\cref{thm:transport}): both winding
matchings carry $\bdelta\ne\bm0$ (the interior matching a bare operator
dipole $\brho\ne\bm0$, the wrapping matching $\brho+\bL\ne\bm0$), while
the genuine contractible process has $\bdelta=\bm0$. Keeping
$\bdelta=\bm0$ retains $\varepsilon$ and deletes \emph{both} winding
matchings, giving the clean $2|\varepsilon|$. Only this scheme reaches
the physical answer.

%=====================================================================
\section{Specialization to QSI, and role in the campaign}
\label{sec:qsi}

For the XXZ pyrochlore the winding loops are the second-order four-site
cycles ($\gfour=4\Jpm^2/\Jzz$) and the wrapping hexagons; the
contractible loops are the genuine $[111]$ hexagons carrying the photon
($\ghex=12|\Jpm|^3/\Jzz^2$). \cref{thm:transport} removes the former at
all orders and keeps the latter at bulk density, so the projected
effective Hamiltonian on the $2970$-state ice manifold of the
$2{\times}2{\times}2$ FCC cluster is a clean, fully diagonalizable model
of the emergent gauge sector.

Because the couplings scale as $\Jpm^k$ exactly, one effective-Hamiltonian
build serves every coupling and both signs of $\Jpm$. This is the
enabling step of the accompanying program: benchmark the projected
$0$-flux ($\Jpm>0$) gauge thermodynamics against sign-problem-free
quantum Monte Carlo, then cross into the $\pi$-flux ($\Jpm<0$) sector --
the QMC dead zone of the cerium pyrochlores -- where the
transport-projected response operators yield the flux-locked
thermodynamic discriminant of the emergent photon.

%=====================================================================
\section{Verdict}
\label{sec:verdict}

Twist averaging, as a means of removing torus-winding contamination from
QSI exact diagonalization, \emph{does not work}: averaging the
Hamiltonian keeps half of every winding coupling
(\cref{cor:half}), and averaging observables keeps even more
(\cref{cor:obs}). Both failures trace to a single misidentification --
assigning a loop's holonomy to an operator that is in fact a sum over
virtual paths (\cref{thm:matched}). The object that succeeds is the
zero-transport projector (\cref{thm:transport}): gauge invariant, exact
at all orders for the winding loops, transparent about its one
higher-order bias, and -- because it acts on processes rather than
spectra -- a cleaning of the entire gauge sector, Hamiltonian and
responses alike. That is what makes a $32$-site pyrochlore a quantitative
laboratory for $\pi$-flux quantum spin ice.

\end{document}
